% !TEX root = p2-report.tex
\documentclass[UTF8,a4paper,zihao=5,fontset=none]{ctexart}
\usepackage[margin=22mm,headheight=15pt,headsep=7mm,footskip=10mm]{geometry}
\usepackage{amsmath,amssymb}
\usepackage{graphicx}
\usepackage{float}
\usepackage{xcolor}
\usepackage{fancyhdr}
\usepackage[unicode,colorlinks=true]{hyperref}

\IfFontExistsTF{SimSun}{
  \setCJKmainfont{SimSun}[BoldFont=SimHei]
  \setCJKsansfont{Microsoft YaHei}
}{
  \IfFontExistsTF{Noto Serif CJK SC}{
    \setCJKmainfont{Noto Serif CJK SC}
    \setCJKsansfont{Noto Sans CJK SC}
  }{
    \setCJKmainfont{FandolSong-Regular}[BoldFont=FandolSong-Bold]
    \setCJKsansfont{FandolHei-Regular}
  }
}
\IfFontExistsTF{Times New Roman}{\setmainfont{Times New Roman}}{\setmainfont{TeX Gyre Termes}}
\definecolor{ink}{HTML}{172326}
\definecolor{muted}{HTML}{627376}
\definecolor{accent}{HTML}{126D70}
\hypersetup{
  pdftitle={文档图像偏色与光照不均：成因、建模与校正},
  pdfsubject={Machine Vision P1},
  pdfauthor={},
  linkcolor=accent,citecolor=accent,urlcolor=accent
}
\color{ink}
\linespread{1.42}
\setlength{\parindent}{2em}
\setlength{\parskip}{0.25em}
\setlength{\emergencystretch}{1em}
\clubpenalty=10000
\widowpenalty=10000
\displaywidowpenalty=10000
\ctexset{
  section={format=\large\bfseries,beforeskip=1em,afterskip=0.5em},
  subsection={format=\normalsize\bfseries,beforeskip=0.8em,afterskip=0.4em}
}
\pagestyle{fancy}
\fancyhf{}
\fancyhead[L]{\small\color{muted}Machine Vision / P1}
\fancyhead[R]{\small\color{muted}文档图像质量技术综述}
\fancyfoot[C]{\small\color{muted}\thepage}
\renewcommand{\headrulewidth}{0.3pt}
\graphicspath{{../figures/}}

\usepackage{booktabs,tabularx,caption}
\IfFontExistsTF{SimSun}{\setCJKmonofont{SimSun}}{\IfFontExistsTF{Noto Sans CJK SC}{\setCJKmonofont{Noto Sans CJK SC}}{\setCJKmonofont{FandolFang-Regular}}}
\captionsetup{font=small,labelfont=bf,skip=5pt}
\hypersetup{pdftitle={板书图像多尺度特征匹配质量评测：可重复性消融实验},pdfsubject={Machine Vision P2 synthetic validation}}
\fancyhead[L]{\small\color{muted}Machine Vision / P2}
\fancyhead[R]{\small\color{muted}多尺度特征匹配消融实验}
\graphicspath{{../figures/p2/}}
\linespread{1.25}
\newcommand{\filepath}[1]{\nolinkurl{#1}}
\newcommand{\R}{\mathcal{R}}
\newcommand{\Sset}{\mathcal{S}}
\newcommand{\Mset}{\mathcal{M}}

\begin{document}
\begin{center}
{\LARGE\bfseries 板书图像多尺度特征匹配质量评测}\par
\vspace{3mm}
{\large 可重复性消融实验}\par
\vspace{2mm}
{\small Lab2 / 第 3 周 / 周四 \qquad 2026 年 9 月 18 日}\par
{\small 合成代理数据验证版}
\end{center}

\begin{abstract}
面向拍摄板书与扫描试卷中的特征不稳定问题，建立单文件评测流程，比较 Sobel、LoG、Harris、Shi-Tomasi 和 SIFT 五类完整匹配管线，考察理想及高斯低通、高通预处理、外部图像金字塔与非极大值抑制的作用。采用共同可见区域内的几何重复率、双向比例检验后的匹配正确率、条件定位误差和每图对运行时间，并绘制图像尺度、附加噪声及金字塔深度曲线。实验对象为十组程序生成的试卷风格图像，几何控制点由生成器给出。结果用于验证评测口径及分析参数敏感性，不作为真实板书效果的证明。对预处理和多尺度的评价应同时考虑空间覆盖、误差、匹配数量与时间，不能仅由一项比例指标判定可信度。当前材料未包含自采照片或人工同名点标注，因而尚未满足作业的数据采集条件。
\end{abstract}

\section{任务范围与数据边界}
作业要求不少于十张不同光照、尺度、模糊条件的自采板书或试卷图像，并为每张场景图相对于参考图手工标注至少二十对同名点。当前仓库没有这类照片，因此本报告只陈述实际完成的合成代理实验，不将生成控制点称为人工标注。

代理集含十组参考/场景图，共二十张 $960\times720$ PNG，每组有三十二对生成控制点。内容包括文字、线框、交叉线与图形，十组共享基本版式但改变文字、颜色和部分细节；场景引入透视扰动、光照变化、尺度、模糊与噪声。原生尺度因子为 $0.86$ 至 $1.06$，高斯模糊参数为 $0$ 至 $2.0$ 像素，原生附加噪声标准差为 $1$ 至 $12$ 灰度级。各因素在这些样本中同时变化，不能由逐图比较单独归因于某一种原生扰动。

数据清单将来源记为 \filepath{synthetic_procedural}，标注状态记为 \filepath{generated_control_points_not_manual}。合成数据使用生成时的场景到参考单应矩阵 $H$，其控制点重投影误差只验证坐标一致性，不代表人工标注质量。实际采集时可共用清晰正视参考图；未提供 $H$ 时，程序以全部人工点拟合单应矩阵，并输出重投影 RMSE。该拟合残差不是独立测试误差，纸面弯曲、镜头畸变或误标都可能使平面模型失效。

\section{方法与单文件实现}
全部 P2 实现位于 \filepath{p2_multiscale_feature_matching.py}，包括代理数据生成、清单校验、特征与描述子、匹配、指标、曲线和自检。流程为灰度读取、各金字塔层缩放、频域预处理、特征提取、坐标回映射与融合、共同可见区域筛选、描述子匹配和几何评估。关键步骤使用中文注释，报告及数据文件不增加第二个 Python 实现。

\subsection{空域响应与描述子}
Sobel 先作 $3\times3$ 高斯平滑，再计算 $3\times3$ 导数核，响应为
\begin{equation}
r_{\mathrm{Sobel}}(x,y)=\sqrt{I_x(x,y)^2+I_y(x,y)^2}.
\end{equation}
该操作强调灰度变化\cite{sobel}，本实验将高响应位置采样为点；它并非与角点检测器完全等价的兴趣点定义。

LoG 以标准差 $1.2$ 的高斯平滑后接 $3\times3$ 离散 Laplacian，取绝对响应：
\begin{equation}
r_{\mathrm{LoG}}(x,y)=\left|\nabla^2(G_{1.2}*I)(x,y)\right|.
\end{equation}
这里实现的是固定层内尺度的平滑后二阶响应采样，而非零交叉边缘检测或尺度归一化 LoG 斑点检测。OpenCV 的 Laplace 教程给出二阶导数形式并演示平滑后计算\cite{laplace}。

角点方法用局部结构张量及其特征值 $\lambda_1,\lambda_2$：
\begin{equation}
M=\sum_{(x,y)\in W}w(x,y)
\begin{bmatrix}I_x^2&I_xI_y\\I_xI_y&I_y^2\end{bmatrix},\qquad
r_H=\det(M)-0.04\,\mathrm{tr}(M)^2,\quad
r_{ST}=\min(\lambda_1,\lambda_2).
\end{equation}
$W$ 为局部窗口，$w$ 为窗口权重。Harris 使用 \filepath{cornerHarris}，邻域宽二、导数核宽三，仅保留正响应；Shi-Tomasi 使用邻域宽三、导数核宽三的 \filepath{cornerMinEigenVal}。两种评分的定义可见官方教程\cite{harris,shi}。不直接调用带内置距离筛选的 \filepath{goodFeaturesToTrack}，以便显式切换共同 NMS。

SIFT 调用 OpenCV \filepath{SIFT_create}，每层候选预算二百二十，参数为 \filepath{contrastThreshold=0.01}、\filepath{edgeThreshold=10}、\filepath{sigma=1.2}。其尺度空间极值、方向分配和原生一百二十八维描述子与官方介绍一致\cite{sift}；本实验使用的是 OpenCV 实现，不是从零重写 SIFT。

其余四种方法取 $32\times32$ 邻域，面积缩放为 $16\times16$，去均值并作 $L_2$ 归一化，得到二百五十六维补丁描述子；近常量补丁被舍弃。该描述子没有 SIFT 的方向与尺度规范化，因此方法比较针对完整管线，不能将匹配正确率差异只归于检测器。

\subsection{频域预处理}
令 $F=\mathrm{fftshift}(\mathrm{FFT2}(I/255))$，滤波后图像为
\begin{equation}
J=\Re\left\{\mathrm{IFFT2}\left[\mathrm{ifftshift}(T\odot F)\right]\right\}.
\end{equation}
$T$ 为传递函数，$\odot$ 为逐元素乘积。该移频、掩模、逆移频与逆变换链路参照官方 Fourier 教程\cite{fourier}。频率半径使用移频后的整数直流中心，奇数宽高也用整除。若当前层宽高为 $w,h$，令
\begin{align}
D(u,v)&=\sqrt{(u-\lfloor w/2\rfloor)^2+(v-\lfloor h/2\rfloor)^2},\\
T_{IL}&=\mathbf{1}[D\leq0.12\min(w,h)],&T_{IH}&=1-T_{IL},\\
T_{GL}&=\exp\left[-\frac{D^2}{2\{0.10\min(w,h)\}^2}\right],&T_{GH}&=1-T_{GL}.
\end{align}
$IL,IH,GL,GH$ 分别表示理想低通、理想高通、高斯低通和高斯高通。另设不滤波对照。截止半径按层尺寸变化，比例为本实验约定，并非普适最优值。低通结果裁剪到 $[0,1]$；高通结果按每图百分之一至百分之九十九分位作线性映射，再转为八位灰度。因此高通比较包含自适应对比度归一化，不是仅替换线性滤波核。未作 FFT 边界填充，周期边界假设也属于实现限制。

\subsection{多尺度融合与 NMS}
外部金字塔深度 $L\in\{1,2,3\}$，层内尺寸比例为 $2^{-l}$，$l=0,\ldots,L-1$；粗层以 \filepath{INTER_AREA} 面积缩放生成，不声称实现了严格的连续高斯尺度空间。每层分别预处理和提取，按实际尺寸比例将坐标映回原图。横坐标为
\begin{equation}
x_0=(x_l+0.5)\frac{w_0}{w_l}-0.5,
\end{equation}
纵坐标同理。半像素项与 OpenCV resize 的像素中心约定一致，避免直接乘二导致粗层定位偏移。四种响应采样方法的各层响应独立归一化，SIFT 保留原生响应分数；融合时乘 $1/(1+0.2l)$，所有层共享最终二百二十点预算，因此增加层数不等于增加最终点数。

四种响应采样方法固定使用第九十九百分位阈值；开启 NMS 时先保留 $5\times5$ 局部最大响应，再对融合点按分数降序实施五原图像素的距离抑制。关闭 NMS 不改变阈值和最终预算，同时关闭这两步。SIFT 的 NMS 消融只关闭共同后处理距离抑制，内部 DoG 极值筛选仍保留；同样，关闭外部金字塔并不关闭 SIFT 的内部尺度空间。

\section{指标定义与实验控制}
令 $\pi_H(q)$ 为场景点 $q$ 经 $H$ 映到参考坐标后的非齐次位置，距离
\begin{equation}
e(p,q)=\|p-\pi_H(q)\|_2
\end{equation}
始终以参考图像素计量。先筛去经双向几何映射落到另一幅图像范围外的点，得到共同可见区域的点集 $\R,\Sset$，避免将不可见裁切位置计入分母。该筛选检查中心点，不保证描述子窗口全部位于共同区域。

\subsection{重复率与定位误差}
枚举 $e(p,q)\leq\epsilon_r$ 的候选点对，按距离从小到大贪心保留一对一关系，得到 $\Mset_r$，其中 $\epsilon_r=5$ 像素。定义
\begin{equation}
\mathrm{Rep}=\frac{|\Mset_r|}{\min(|\R|,|\Sset|)},\qquad
E_{\mathrm{loc}}=\mathrm{median}_{(p,q)\in\Mset_r}e(p,q).
\end{equation}
分母为零时重复率记零；无重复点时定位误差记为未定义。贪心配对不是最大基数二分匹配，也不是全局最优分配；该口径在所有设置中固定使用。定位误差仅针对五像素以内被重复检测的点，不是全部匹配或全部标注点的定位精度，低误差必须结合覆盖与重复点数解读。

\subsection{匹配正确率与时间}
描述子以 $L_2$ 最近邻匹配；正反方向均要求最近邻/次近邻距离比小于 $0.78$，且互为最近邻，得到接受集 $\Mset_a$。不使用被测匹配估计 $H$ 再自我判断正确性。定义
\begin{equation}
\mathrm{Prec}=\frac{\sum_{(p,q)\in\Mset_a}\mathbf{1}[e(p,q)\leq4]}{|\Mset_a|},\qquad
t=1000(t_{\mathrm{end}}-t_{\mathrm{start}})\ \mathrm{ms}.
\end{equation}
无接受匹配时正确率未定义而不是零，CSV 中为 \filepath{nan}、JSON 中为 \filepath{null}。汇总只平均有限值，并附有效样本数；这不是把全部匹配池化后的总体正确率。逐图文件同时保留接受匹配数、正确匹配数、几何重复数、原始及可见区域点数和匹配误差，便于发现少量点导致的高比例。

计时覆盖两幅图的预处理、检测、描述、几何评估和匹配，不含磁盘 I/O；每个图对/配置只计时一次，没有预热或多次重复，跨十图的标准差是样本差异，不是计时稳定性置信区间。OpenCV 固定单线程。

\subsection{实验矩阵与可重复性}
主矩阵为十图对、五方法、五预处理与三个外部深度，共七百五十行。消融以高斯低通、三层、NMS 开启为 \filepath{full}，分别关闭外部金字塔、频域滤波或共同 NMS，共二百行。尺度实验固定三层与 NMS，仅将场景图额外缩放至 $0.5,0.75,1.0,1.25$ 倍，参考不变，并同步组合逆 resize 与 $H$，共一千行。尺度横轴表示数字重采样倍率，不是实测焦距或拍摄距离。

噪声实验对场景图额外加入标准差 $0,5,10,20$ 灰度级的高斯噪声，原生扰动仍保留，共一千行。参考图不加噪；所有方法及滤波共用同一样本的噪声实例，同一标准随机场仅改变振幅，随后裁剪到八位范围。随机种子固定为 $20260917$。曲线均值及误差棒按十图汇总，误差棒是样本标准差，不是置信区间。

本次环境为 Python 3.12.14、OpenCV 4.14.0、NumPy 2.5.3、Matplotlib 3.11.2、Pillow 12.3.0；精确发行包版本见 \filepath{requirements-p2.txt}。\filepath{summary.json} 记录参数、平台、线程数、完成时间以及代码、图像和标注 SHA-256。内置自检覆盖低通直流保持、高通去直流、NMS、共同可见区域、一对一配对、数值排序、无匹配口径及噪声确定性等行为。

\clearpage
\section{结果与曲线}
\subsection{预处理对比及完整设置}
表~\ref{tab:filters} 固定三层和 NMS，给出各预处理的十图平均重复率。该条件下高斯低通的 Harris、Shi-Tomasi、SIFT 分别为 $0.807,0.792,0.799$；而 Sobel 在理想低通下高于高斯低通。LoG 在理想低通和两种高通下的重复率均高于高斯低通。这些结果表明预处理与管线有交互作用，不能将同一滤波判为对所有方法统一最优。

\begin{table}[H]
\centering\small
\caption{三层、NMS 开启时的重复率。每格为十图均值，数值越高表示本口径下重复点比例越大；结果来自七百五十行主矩阵中的三层子集。}
\label{tab:filters}
\setlength{\tabcolsep}{7pt}
\begin{tabular}{lrrrrr}
\toprule
方法 & 不滤波 & 理想低通 & 高斯低通 & 理想高通 & 高斯高通\\
\midrule
Sobel & 0.213 & 0.567 & 0.397 & 0.474 & 0.400\\
LoG & 0.382 & 0.560 & 0.508 & 0.545 & 0.551\\
Harris & 0.680 & 0.661 & 0.807 & 0.658 & 0.610\\
Shi-Tomasi & 0.747 & 0.673 & 0.792 & 0.691 & 0.688\\
SIFT & 0.552 & 0.643 & 0.799 & 0.548 & 0.551\\
\bottomrule
\end{tabular}
\end{table}

表~\ref{tab:full} 对应消融的完整设置，不是为每种方法逐一调出的最好配置。SIFT 在该设置下的条件定位误差最低，但运行时间高于其他四条管线；Harris 的重复率及正确率较高。LoG 的重复率 $0.508$ 并不能保证描述子匹配可信，其平均正确率仅为 $0.723$，平均接受匹配数为 $8.8$、正确匹配数为 $6.9$。比例和数量应同时查看，不能把本代理集的结果当作普遍的检测器排名。

\begin{table}[H]
\centering\small
\caption{完整设置的四项指标。高斯低通、三层、NMS 开启；十图均值，正确率有效样本均为十。误差列为每图重复点中位误差的跨图平均，时间为消融运行阶段的单次计时均值。}
\label{tab:full}
\begin{tabular}{lrrrr}
\toprule
方法 & 重复率 & 匹配正确率 & 定位误差 / px & 时间 / ms\\
\midrule
Sobel & 0.397 & 0.744 & 1.313 & 342.4\\
LoG & 0.508 & 0.723 & 1.545 & 363.2\\
Harris & 0.807 & 0.972 & 0.960 & 348.1\\
Shi-Tomasi & 0.792 & 0.971 & 1.003 & 327.6\\
SIFT & 0.799 & 0.912 & 0.413 & 622.3\\
\bottomrule
\end{tabular}
\end{table}

\begin{figure}[H]
\centering
\includegraphics[width=\textwidth]{sample_matches.png}
\caption{几何正确的匹配和少量错误匹配可以同时出现。第一组合成图的 SIFT 完整设置示例，左为参考、右为场景；仅画前四十个接受匹配，绿线误差不超过四参考像素，红线超过阈值。示例不能代替十图统计，也不是自采图。}
\end{figure}

\clearpage
\section{组件消融与敏感性分析}
表~\ref{tab:delta} 给出各关闭设置相对于完整设置的变化量。关闭外部金字塔使 Sobel、Harris、Shi-Tomasi、SIFT 的重复率降低，但 LoG 上升约四个百分点；同时 LoG 正确率上升约二十一个百分点。这与固定预算下的层候选竞争及不同尺度补丁支持域有关的解释相容，但当前结果不独立证明这一机制。

关闭频域滤波时，Harris 和 SIFT 的重复率分别下降 $12.73$ 和 $24.64$ 个百分点，而 LoG 的正确率上升 $27.37$ 个百分点，平均正确匹配数由 $6.9$ 增到 $58.7$。因此高斯低通对该 LoG 管线在几何重复率和描述子可信度上产生不同方向的变化。其他四种方法的频域去除均降低重复率，不能由这一现象推出所有实际照片都应启用低通。

\begin{table}[H]
\centering\small
\caption{消融变化量，单位为百分点，等于关闭设置均值减去完整设置均值；负值表示相应比例下降。P、F、N 分别表示关闭外部金字塔、频域滤波、共同 NMS。各设置保持响应阈值和最终特征预算不变。}
\label{tab:delta}
\setlength{\tabcolsep}{6pt}
\begin{tabular}{lrrrrrr}
\toprule
 & \multicolumn{3}{c}{重复率变化 / pp} & \multicolumn{3}{c}{正确率变化 / pp}\\
\cmidrule(lr){2-4}\cmidrule(lr){5-7}
方法 & P & F & N & P & F & N\\
\midrule
Sobel & $-6.01$ & $-18.45$ & $-29.41$ & $+0.91$ & $+5.48$ & $+5.31$\\
LoG & $+4.01$ & $-12.59$ & $-35.18$ & $+21.04$ & $+27.37$ & $-27.54$\\
Harris & $-1.55$ & $-12.73$ & $-22.18$ & $+0.41$ & $-0.02$ & $+0.96$\\
Shi-Tomasi & $-0.23$ & $-4.50$ & $-11.00$ & $+1.02$ & $+0.32$ & $+2.82$\\
SIFT & $-5.11$ & $-24.64$ & $-14.27$ & $+0.43$ & $-7.07$ & $+4.55$\\
\bottomrule
\end{tabular}
\end{table}

关闭 NMS 时所有方法的重复率下降，但部分正确率和条件定位误差改善。Harris 的定位误差由 $0.960$ 降至 $0.585$ 像素，平均正确匹配数却由 $117.3$ 降至 $50.0$；SIFT 的正确率由 $0.912$ 升至 $0.958$，正确匹配数由 $88.2$ 降至 $72.9$。这说明更小的条件误差或更高的比例不等价于更高的覆盖与总匹配能力。无 NMS 的 Sobel、LoG 正确率只各有七、八个有效图对，无外部金字塔的 LoG 有九个；这些均值的样本集合不同，差值应谨慎解读。

\clearpage
\begin{figure}[H]
\centering
\includegraphics[width=\textwidth]{ablation_metrics.pdf}
\caption{关闭组件并不使所有指标同方向变化。每种方法从左至右依次为 full、关闭外部金字塔、关闭频域、关闭共同 NMS；full 为高斯低通、三层、NMS 开启。误差棒是跨图样本标准差，不是置信区间；重复率和正确率坐标范围为零至 $1.02$，越界标准差部分不显示。定位误差针对几何重复点，时间含两图完整评测、不含 I/O。}
\end{figure}
图中运行时间也不是统一方向的收益证明。关闭频域使四种响应采样管线明显减少本次均值时间；SIFT 本次从 $622.3$ 变为 $691.9$ ms，反而上升。滤波会改变被 SIFT 内部处理的图像内容，可能影响候选数量和内部计算，但没有阶段计时证据，不能据此确认原因；单次测量及系统负载也限制了精确速度结论。

\clearpage
\subsection{图像数字尺度敏感性}
高斯低通下，场景倍率从 $1.0$ 缩为 $0.5$ 时，Harris 重复率由 $0.807$ 降至 $0.521$，SIFT 由 $0.799$ 降至 $0.610$；SIFT 放大到 $1.25$ 时为 $0.778$。但 Sobel 在 $0.5$ 倍处为 $0.560$，高于 $1.0$ 倍的 $0.397$；其平均重复点数却由 $87.4$ 减少到 $72.7$，平均场景可见点数由 $220.0$ 减少到 $128.9$。分母和覆盖发生变化，比例上升不代表重复点数量增加。重采样、滤波支持域及候选竞争也共同影响曲线，不能把该现象视为真实拍摄缩小的普遍收益。横轴只改变场景数字尺寸，不改变参考图及误差单位。

\begin{figure}[H]
\centering
\includegraphics[width=\textwidth]{repeatability_scale.pdf}
\caption{图像尺度与预处理共同影响几何重复率。横轴为场景额外数字缩放倍率，参考图不变，三层和 NMS 固定开启。none 为不滤波，LP/HP 为低通/高通；每面板含五方法十图均值和样本标准差，以颜色、标记和线型区分。误差仍以参考图像素判断，比例坐标限定零至 $1.02$，越界标准差部分不显示；不代表真实拍摄距离曲线。}
\end{figure}

\clearpage
\subsection{附加噪声敏感性}
附加噪声从零增至二十 DN 时，不滤波 Harris 的重复率由 $0.680$ 降至 $0.524$，SIFT 由 $0.552$ 降至 $0.347$；高斯低通下分别由 $0.807$ 到 $0.797$、$0.799$ 到 $0.756$。在当前参数和代理集上，两者的低通几何重复率较稳定；轻微非单调波动不应解读为噪声带来收益。理想低通 SIFT 末端为 $0.610$，低于高斯低通；两种高通 SIFT 末端约 $0.342,0.350$。这些曲线只测几何重复性，不能推翻 LoG 消融所显示的描述子正确率与数量限制。

\begin{figure}[H]
\centering
\includegraphics[width=\textwidth]{repeatability_noise.pdf}
\caption{重复率随附加噪声变化的趋势依赖方法及预处理。横轴为场景额外高斯噪声标准差，单位 DN，即八位灰度级；三层与 NMS 固定，原生噪声和其他扰动仍保留。none、LP、HP 含义同尺度图；点为十图均值，误差棒为样本标准差，非置信区间。坐标限定零至 $1.02$，越界标准差部分不显示；零附加噪声并不意味着原图无噪声。}
\end{figure}

\clearpage
\subsection{外部金字塔深度}
高斯低通下，Sobel 的重复率从一层 $0.337$ 到二层 $0.399$，三层 $0.397$；Harris 从 $0.791$ 到 $0.805$、$0.807$；Shi-Tomasi 从 $0.790$ 到 $0.791$、$0.792$。SIFT 从 $0.748$ 到 $0.793$、$0.799$，LoG 则从 $0.548$ 降至 $0.517$、$0.508$。新增层的收益并不单调或普遍，两层到三层的变化还应结合标准差和额外计算代价判断。
\begin{figure}[H]
\centering
\includegraphics[width=\textwidth]{repeatability_pyramid.pdf}
\caption{增加外部金字塔深度不保证重复率持续提升。横轴是层数，一层表示仅原图、二三层分别加入 $1/2$ 与 $1/4$ 图像；最终点预算仍为二百二十，NMS 开启。五面板分别对应五预处理，点和误差棒为十图均值与样本标准差；比例坐标范围为零至 $1.02$，越界标准差部分不显示。SIFT 各设置均保留内部尺度空间，不能将该横轴解释为 SIFT 总尺度层数。}
\end{figure}


\clearpage
\section{计算摄影链路与可信度讨论}
板书或试卷的匹配结果并非只由末端检测器决定。可将关注链路写为“内容与照明、成像光学与运动、采样和图像处理、频域/多尺度预处理、特征与描述子、匹配和几何判定”。本实验只模拟部分后端图像变化，没有采集 RAW、标定真实噪声、复现相机 ISP 或测量下游 OCR；以下机制解释是与算法及代理现象一致的工作假设，不是由当前消融证明的物理因果关系。

照明的缓慢变化与细笔画、边缘的快速变化占据不同频率范围，但真实信号与噪声并不按频率完全分离。低通可能抑制高频扰动，也会牺牲细笔画及局部对比度；高通可能削弱平缓亮度变化，同时增强噪声并改变描述子的灰度分布。官方 Fourier 教程讨论了图像边缘、噪声的高频成分及突变掩模相关振铃\cite{fourier}。本实现使用圆形理想截止，不能把教程中的矩形掩模现象当作本数据上已量化的振铃证据。

多尺度的作用一方面是以粗层降低局部扰动，另一方面是改变有效邻域与特征预算分配。对补丁描述子，粗层采样意味着更大的原图支持域；对 SIFT，外部层叠加在已存在的内部尺度空间之外。若新增粗层候选挤占原图层稳定位置，或粗层细笔画已丢失，则重复率不一定上升。当前的响应归一化和层权重也是这一融合策略的一部分，不能把观察到的变化解释为所有图像金字塔的共同规律。

NMS 的价值不只在于改变点数，而在于减少近邻重复响应、改善空间覆盖及描述子唯一性。关闭它可能在同一笔画附近产生很多几何上接近的点，影响比例检验的最近/次近邻关系。因此可信度至少应联合检查重复率、定位误差、接受和正确匹配数量、分布覆盖及运行时间。即使正确率很高，少量或集中在一处的匹配仍不足以支持稳定几何估计。本实验没有估计置信概率，正确率也不是某一个具体匹配的概率校准值。

\section{局限与后续真实采集}
首先，十图共享程序版式，且仅覆盖已设定的扰动；生成控制点与已知 $H$ 给出理想几何参照，无法代替真实拍摄中的畸变、反光、压缩、复杂笔迹或人工误标。因此结论只适用于这套代理数据和参数，不宣称某方法普遍优于其他方法。

其次，SIFT 与其他方法的描述子不同，且关闭外部 NMS/金字塔时内部机制仍在；当前实验回答完整管线在共同预算下如何变化，不是检测器的纯净因果比较。若需要更严格的检测器消融，可另作共同描述子评测，但那属于新的实验而非本报告已经完成的结果。

第三，重复率采用贪心配对，定位误差是条件统计，运行时间仅一次测量。正式分析还应复测计时、查看逐图差异，并考虑以全局一对一分配作敏感性检查；当前不报告显著性检验或独立原稿上的泛化收益。

满足课程数据条件的直接步骤是实际采集不少于十张场景照片，按真实情况记录条件，与清晰参考图逐对手工标注至少二十点并复核。仓库提供 \filepath{real-manifest.template.json} 和含二十个空点对的 \filepath{manual-annotation.template.json}，空值不能通过校验。仅在确实完成采集和复核后标记 \filepath{self_captured} 与 \filepath{manual_verified}。程序的严格校验可查字段和坐标，不可能鉴定数据来源真实性。真实结果应写入独立输出目录，并替换报告结果与结论，不能沿用本报告的合成数值。

\section{结论与复现步骤}
本次完成了单文件 P2 评测实现、三项单因素消融、尺度与噪声曲线，以及逐图指标和可核对的运行记录。实验支持的结论是：滤波、尺度融合和 NMS 需要按完整匹配管线及共同预算评价，关闭某一组件不必然降低所有指标；可信度不能简化为重复率、正确率或定位误差中的单一数值。自采和人工标注尚未完成，所以本材料是可复现的代理验证版，而非已满足全部数据要求的最终作业。

\noindent\begin{minipage}{\linewidth}
使用 Python 3.12 创建独立环境，在仓库或提交包根目录运行：
\begin{verbatim}
python -m pip install -r requirements-p2.txt
python p2_multiscale_feature_matching.py --self-test
python p2_multiscale_feature_matching.py --validate-only
python p2_multiscale_feature_matching.py --data-root data/p2
\end{verbatim}
\end{minipage}
已有 PNG 是严格复现优先输入，重新生成会受系统字体影响。切换真实数据时先校验，保留代理结果：
\begin{verbatim}
python p2_multiscale_feature_matching.py \
  --data-root data/private/p2 --require-manual --validate-only
python p2_multiscale_feature_matching.py \
  --data-root data/private/p2 --require-manual \
  --output-root output/p2-real --figure-root figures/p2-real
\end{verbatim}
上面的续行反斜杠适用于 POSIX shell；PowerShell 应将各命令写成一行。输出包含 \filepath{main_results.csv}、\filepath{ablation_results.csv}、\filepath{scale_results.csv}、\filepath{noise_results.csv}，以及汇总和 PDF/SVG/PNG 图表。报告不参与数值计算；重新运行后应根据新 CSV 同步更新报告。提交包仅含一个 \filepath{.py}，排除 P1 实现、临时运行库和 Git 元数据。

\begin{thebibliography}{9}
\small
\bibitem{sobel} Ana Huam\'{a}n（原作者），OpenCV contributors（维护）. Sobel Derivatives. OpenCV 官方教程源码，访问日期 2026-09-18。\href{https://github.com/opencv/opencv/blob/cf66d2c6add2e94396eabcaf5e85a4578dc3636b/doc/tutorials/imgproc/imgtrans/sobel_derivatives/sobel_derivatives.markdown}{固定版本源码}。
\bibitem{laplace} Ana Huam\'{a}n（原作者），OpenCV contributors（维护）. Laplace Operator. OpenCV 官方教程源码，访问日期 2026-09-18。\href{https://github.com/opencv/opencv/blob/cf66d2c6add2e94396eabcaf5e85a4578dc3636b/doc/tutorials/imgproc/imgtrans/laplace_operator/laplace_operator.markdown}{固定版本源码}。
\bibitem{harris} OpenCV contributors. Harris Corner Detection. OpenCV 官方教程源码，访问日期 2026-09-18。\href{https://github.com/opencv/opencv/blob/cf66d2c6add2e94396eabcaf5e85a4578dc3636b/doc/py_tutorials/py_feature2d/py_features_harris/py_features_harris.markdown}{固定版本源码}。
\bibitem{shi} OpenCV contributors. Shi-Tomasi Corner Detector \& Good Features to Track. OpenCV 官方教程源码，访问日期 2026-09-18。\href{https://github.com/opencv/opencv/blob/cf66d2c6add2e94396eabcaf5e85a4578dc3636b/doc/py_tutorials/py_feature2d/py_shi_tomasi/py_shi_tomasi.markdown}{固定版本源码}。
\bibitem{sift} OpenCV contributors. Introduction to SIFT (Scale-Invariant Feature Transform). OpenCV 官方教程源码，访问日期 2026-09-18。\href{https://github.com/opencv/opencv/blob/cf66d2c6add2e94396eabcaf5e85a4578dc3636b/doc/py_tutorials/py_feature2d/py_sift_intro/py_sift_intro.markdown}{固定版本源码}。
\bibitem{fourier} OpenCV contributors. Fourier Transform. OpenCV 官方教程源码，访问日期 2026-09-18。\href{https://github.com/opencv/opencv/blob/cf66d2c6add2e94396eabcaf5e85a4578dc3636b/doc/py_tutorials/py_imgproc/py_transforms/py_fourier_transform/py_fourier_transform.markdown}{固定版本源码}。
\end{thebibliography}
\end{document}
